\section{相关工作与理论背景}

本文关心的不是大词汇量识别，而是短语音里一个目标首辅音是否真的出现。这个问题更接近传统语音处理里的短时分析、局部谱形比较和模板匹配，而不是把整段音频直接交给一个黑盒分类器。

\subsection{短时分析与频谱表示}

Davis 和 Mermelstein 早就说明，语音识别里最有用的前端不是整段波形本身，而是经过短时分帧后的参数化表示\cite{davis1980}。做法很简单：把波形切成几十毫秒一帧，给每帧加窗，再做 Fourier transform，最后看这一小段时间里能量落在哪些频率上。这样得到的谱图，既保留了辅音释放、摩擦和元音过渡这些局部线索，也比原始波形稳定得多。

Rabiner 对 HMM 的综述则把这种帧级表示和序列建模连了起来\cite{rabiner1989}。它的核心意思是，语音不是静态图片，真正有意义的是“某个短时结构在时间上怎么移动”。虽然本文没有上 HMM，但这个观点很重要，因为它提醒我们：如果目标只是在整词里检出词首辅音，先把搜索范围压到短时块，再比较局部谱形，通常比直接看整词平均特征更合适。

\subsection{辅音线索为什么更依赖局部谱形}

对 /g/、/b/、/d/、/z/ 这类首辅音来说，最关键的证据往往只出现在很短的一段起始区间。爆破音的释放段短、转移快，摩擦音则在高频段拖得更久；这些差异一旦被整词平均掉，信息就会变少。Wilkinson 和 Russell 在 TIMIT 上做 phone recognition 时，也体现出同一个结论：越接近发音边界的局部线索，越能帮助辅音判别\cite{wilkinson2002}。

因此，本文的主方法没有先上复杂模型，而是回到课堂里最直接的工具链：短时 FFT、局部频谱、归一化相关。这样做的好处不是“更先进”，而是把问题本身讲清楚。

\subsection{本文为什么先用课堂方法}

本文后面的 AI 对照只负责补充结果，不负责定义主线。主线先从局部频谱模板出发，把“一个词首辅音在声学上长什么样”讲明白，再看更复杂的方法能把这个边界推到哪里。这样安排，后面的比较才不会变成单纯的分数表。
